
\begin{table}[ht]
\centering
\footnotesize
\setlength{\tabcolsep}{4pt}
\begin{tabular}{lllcc}
\toprule
\makecell{\textbf{Family}\\\textbf{(mod 256)}} &
\textbf{Example program} & 
\textbf{Its output} &
\makecell{\textbf{Earliest}\\\textbf{round}} &
\makecell{$\mathbb{E}[\text{first round}]$\\\textbf{(univ.\ prior)}} \\
\midrule
Arithmetic & \texttt{S+[.++]}            & $1,3,5,7,9,\ldots$    & 0   & $\approx 105$ \\
Fibonacci  & \texttt{S,[[.C>.C>]}        & $1,1,2,3,5,\ldots$    & 512 & $>53{,}000$ \\
Geometric  & \texttt{S+[.L>]}            & $1,3,9,27,81,\ldots$  & 256 & $>53{,}000$ \\
Quadratic  & \texttt{S,.[<C>{}>VX<RX++]} & $9,25,59,111,\ldots$  & 512 & $>53{,}000$ \\
Cubic      & \texttt{S+[[-.L>L>-]-]}     & $0,254,236,74,\ldots$ & 512 & $>53{,}000$ \\
\bottomrule
\end{tabular}
\vspace{6pt}
\caption{\textbf{Program families with recognizable mathematical structure discovered by the generator during self-play.}  
\emph{Earliest round} gives the earliest round in which a member of the family first appears during training, while \emph{univ. prior} gives the expected first appearance round if programs are drawn from the universal prior, including the added primitives. See \cref{apx:emergent-math-structure} for additional details.}
\label{tab: discovered-math-structure}
\end{table}
